\section{Функция на разпределение}

Функцията на разпределение е едно от основните понятия в теорията на вероятностите. На английски език терминът е \emph{Cumulative Distribution Function} (често съкращавано като CDF).

\begin{keydefn}[функция на разпределение]\label{def:cdf}
Нека $X$ е случайна величина, дефинирана върху вероятностно пространство
$(\Omega, \mathcal{A}, \mathbb{P})$. \emph{Функция на разпределение} на $X$ наричаме
\[ F_X(x) = \mathbb{P}(X < x), \qquad x \in \mathbb{R} , \]
тоест вероятността $X$ да приеме стойност, строго по-малка от $x$.
\end{keydefn}

\emph{Бележка относно терминологията:} В англоезичната литература (а и не само) функцията на разпределение по принцип се дефинира с нестрого неравенство, т.е. $F(x) = \mathbb{P}(X \le x)$. Въпреки това в част от българската традиция — и в този курс — се използва строго неравенство ($X < x$); срещат се и български учебници с нестрогия вариант. На практика това не води до съществени промени, особено при непрекъснати случайни величини, разгледани по-късно. Тази терминологична особеност следва да се има предвид.

\begin{prop}[Свойства на функцията на разпределение]\label{prop:cdf-props}
За всяка случайна величина $X$ функцията $F_X$ притежава следните свойства:
\begin{enumerate}
    \item[а)] $0 \le F_X(x) \le 1$ за всяко $x$;
    \item[б)] $F_X$ е монотонно ненамаляваща: ако $x_1 < x_2$, то $F_X(x_1) \le F_X(x_2)$;
    \item[в)] $\displaystyle \lim_{x \to -\infty} F_X(x) = 0$ и $\displaystyle \lim_{x \to +\infty} F_X(x) = 1$;
    \item[г)] $F_X$ е непрекъсната отляво (при конвенцията със строго неравенство).
\end{enumerate}
\end{prop}

\begin{proof}
а) следва от това, че $F_X(x)$ е вероятност. За б): ако $x_1 < x_2$, то
$\{X < x_1\} \subseteq \{X < x_2\}$, а вероятността е монотонна
(твърдение~\crosslecture{prop:monotone}{за монотонност от лекция 2}). За в): събитията
$\{X < x\}$ намаляват към $\emptyset$ при $x \to -\infty$ и растат към $\Omega$ при
$x \to +\infty$; остава да приложим непрекъснатостта на вероятностната мярка. Свойство г)
се проверява по същия начин чрез растяща редица $x_n \uparrow x$.
\end{proof}

Точките, в които $F_X$ има скок, са точно тези $x$, за които $\mathbb{P}(X = x) > 0$;
големината на скока е самата $\mathbb{P}(X = x)$. Това ще ни трябва по-нататък при
сходимостта по разпределение, където се изискват точки на непрекъснатост.

\subsection{Връзка с разпределението на дискретна случайна величина}

Дискретните случайни величини, разгледани в лекция 4, могат да бъдат описани чрез техните таблици на разпределение. Функцията на разпределение обаче е много по-общ и често по-удобен инструмент. Тя капсулира цялата информация за вероятностите, свързани със случайната величина $X$.

Нека $X$ е дискретна случайна величина, която приема стойности $x_1 < x_2 < x_3 < \dots$ със съответни вероятности $p_1, p_2, p_3, \dots$. (Допускаме, че стойностите са подредени в нарастващ ред).
Как бихме изчислили функцията на разпределение за тази случайна величина?
Тя има следния вид:
\[
F(x) =
\begin{cases}
0, & \text{ако } x \le x_1 \\
p_1, & \text{ако } x \in (x_1, x_2] \\
p_1 + p_2, & \text{ако } x \in (x_2, x_3] \\
p_1 + p_2 + p_3, & \text{ако } x \in (x_3, x_4] \\
\vdots & \vdots
\end{cases}
\]

Защо е така? Ако $x \le x_1$, тогава вероятността $X$ да приеме стойност, строго по-малка от $x_1$, е 0, тъй като $x_1$ е най-малката възможна стойност. Когато $x$ попадне в интервала $(x_1, x_2]$, единствената възможна стойност на $X$, която е строго по-малка от $x$, е $x_1$. Следователно вероятността е $p_1$. Този процес продължава стъпаловидно. Вижда се, че от тази функция лесно можем да възстановим оригиналната таблица с вероятностите, и обратно. Функцията на разпределение дава единен начин за кодиране на информацията на всяка една случайна величина, независимо дали е дискретна или непрекъсната.

\begin{example}[Индикаторна функция]\label{ex:05-1}
Нека $X$ е индикатор на дадено събитие $A$, тоест $X = \ind_A$. Това е най-простата дискретна случайна величина. Тя приема стойност 1 с вероятност $p_1 = \mathbb{P}(A)$ и стойност 0 с вероятност $p_0 = 1 - p_1$.
Графиката на функцията на разпределение $F(x)$ за този индикатор е стъпаловидна. Тя е 0 за всички стойности до 0 включително. В интервала $(0, 1]$ функцията скача на ниво $1 - p_1$, тъй като там $X$ може да бъде само 0. За стойности, по-големи от 1, функцията скача на ниво 1, защото вероятността $X$ да е по-малко от число, по-голямо от 1, е сигурно събитие (тъй като $X$ е или 0, или 1).
\[
F(x) =
\begin{cases}
0, & \text{ако } x \le 0 \\
1 - p_1, & \text{ако } x \in (0, 1] \\
1, & \text{ако } x > 1
\end{cases}
\]

\begin{lecfigure}[Функцията на разпределение на индикатор $\ind_A$ при
$p_1 = \mathbb{P}(A)$. Стъпаловидна е, с два скока — в $0$ и в $1$. Плътната точка
показва коя стойност се приема в самата точка на скока: при конвенцията
$F(x) = \mathbb{P}(X < x)$ функцията е непрекъсната отляво.\label{fig:cdf-indicator}]
\begin{tikzpicture}[x=2.1cm, y=2.3cm]
  \draw[axisline] (-0.75,0) -- (2.3,0) node[right, font=\small] {$x$};
  \draw[axisline] (0,-0.12) -- (0,1.32) node[above, font=\small] {$F(x)$};
  % level guides
  \draw[guide] (0,1) -- (2.1,1);
  \draw[guide] (0,0.62) -- (1,0.62);
  \node[left, font=\small] at (-0.03,1) {$1$};
  \node[left, font=\small] at (-0.03,0.62) {$1-p_1$};
  \node[below, font=\small] at (1,-0.04) {$1$};
  % the three pieces
  \draw[cdfstep] (-0.7,0) -- (0,0);
  \draw[cdfstep] (0,0.62) -- (1,0.62);
  \draw[cdfstep] (1,1) -- (2.1,1);
  % closed / open endpoints
  \fill[defframe] (0,0) circle (1.7pt);
  \draw[defframe, fill=white, line width=0.7pt] (0,0.62) circle (1.7pt);
  \fill[defframe] (1,0.62) circle (1.7pt);
  \draw[defframe, fill=white, line width=0.7pt] (1,1) circle (1.7pt);
\end{tikzpicture}
\end{lecfigure}

\end{example}

\section{Математическо очакване}

Следващата основна тема е свързана с две изключително важни числови характеристики на случайните величини: \textbf{математическо очакване} (на английски \emph{expectation}) и \textbf{дисперсия} (на английски \emph{variance} или \emph{variation}).

\subsection{Дефиниция на математическо очакване}

Първо ще разгледаме формалната дефиниция за дискретна случайна величина.

\begin{keydefn}
Нека $X$ е дискретна случайна величина, приемаща стойности $x_1, x_2, \dots$ със съответни вероятности $p_1, p_2, \dots$. Тогава под \emph{математическо очакване} на $X$, което бележим с $\E X$, разбираме сумата:
\[ \E X = \sum_{j} x_j p_j \]
ако редът от абсолютните стойности е сходящ, т.е.:
\[ \sum_{j} |x_j| p_j < \infty \]
\end{keydefn}

\textbf{Коментари към дефиницията:}
\begin{enumerate}
    \item Ако $X$ приема \emph{краен брой стойности} $\{x_1, \dots, x_n\}$, то сумата винаги е крайна и математическото очакване винаги съществува. То се задава просто като $\E X = \sum_{j=1}^n x_j p_j$.
    \item Ако $X$ приема \emph{безкраен брой стойности}, е възможно очакването да не съществува. Например, ако разгледаме случайна величина $X$, която приема стойности $j = 1, 2, 3, \dots$ с вероятности $p_j = \frac{6}{\pi^2 j^2}$, тогава очакването би било сума от вида $\sum_{j} j \cdot \frac{6}{\pi^2 j^2} = \frac{6}{\pi^2} \sum_{j} \frac{1}{j}$. Това е хармоничният ред, който е разходящ (сумата му е безкрайност), следователно очакването за тази случайна величина не съществува.
\end{enumerate}

\subsection{Интерпретация на математическото очакване}

\textbf{Физическа интерпретация:} Математическото очакване има много ясен физически смисъл. То играе ролята на \emph{център на тежестта}. Представете си, че имате единична маса, която е разпределена в различни точки (локации) $x_j$ по числовата ос, и всяка точка има маса $p_j$. Тогава очакването $\E X$ е точно координатата на центъра на тежестта на тази система. Както в механиката центърът на тежестта е точката, в която обектът се уравновесява, така и очакването е „центърът“ на разпределението на вероятностите.

\textbf{Математическа интерпретация (Минимизиране на грешката):} Представете си, че трябва да изберете само едно реално число $a$, което най-добре да характеризира дадена случайна величина $X$. Какво означава „най-добре“? Ако искаме да минимизираме \emph{средноквадратичната грешка}, тоест търсим числото $a$, което минимизира израза:
\[ \sum_{j} (x_j - a)^2 p_j \]
се оказва, че този минимум се достига точно когато $a = \E X$. В средноквадратичен смисъл, очакването най-добре описва „типичната“ стойност на случайната величина. Това свойство е фундаментално и ще бъде ключово по-нататък, когато се изучава условното математическо очакване.

\subsection{Примери за математическо очакване}

\begin{example}[Равномерно разпределение]\label{ex:05-2}
Нека $X$ е равномерно разпределена върху числата $\{1, 2, \dots, n\}$. Тоест, вероятността да приеме всяка една от тези стойности е $1/n$ ($p_j = 1/n$ за $j=1,\dots,n$).
Тогава очакването е:
\[ \E X = \sum_{j=1}^n x_j p_j = \sum_{j=1}^n j \frac{1}{n} = \frac{1}{n} \sum_{j=1}^n j \]
Това е просто средноаритметичното на числата.

\end{example}

\begin{example}[Две стратегии на рулетка]\label{ex:05-3}
Нека разгледаме игра на европейска рулетка, която има 37 числа (от 0 до 36).
\begin{itemize}
    \item \emph{Стратегия 1 (Залагане на цвят):} Залагаме на „черно“. Има 18 черни, 18 червени числа и една зелена нула. Залагаме 1 лев. Печелим 1 лев с вероятност $18/37$ и губим 1 лев (тоест печелим $-1$) с вероятност $19/37$. Нека $X$ е печалбата от тази стратегия.
    \[ \E X = (-1) \cdot \frac{19}{37} + 1 \cdot \frac{18}{37} = -\frac{1}{37} \]

    \item \emph{Стратегия 2 (Залагане на конкретно число):} Залагаме 1 лев на числото „1“. Печелим 35 лева чисто (връщат ни 36 лева) с вероятност $1/37$. Губим 1 лев с вероятност $36/37$. Нека $Y$ е печалбата от тази стратегия.
    \[ \E Y = (-1) \cdot \frac{36}{37} + 35 \cdot \frac{1}{37} = \frac{-36 + 35}{37} = -\frac{1}{37} \]
\end{itemize}

От гледна точка на математическото очакване, двете игри са напълно еквивалентни. Средно, дългосрочно, при всеки заложен лев вие губите $1/37$ от лева (около $2{,}7$ стотинки), независимо коя от двете стратегии ще изберете. Това показва, че очакването е мярка за средния резултат. Въпреки това, интуицията ни подсказва, че двете игри се усещат много различно. Тази разлика ще бъде обяснена и количествено измерена, когато въведем понятието \emph{дисперсия}.

\end{example}

\subsection{Свойства на математическото очакване}

Преди да продължим, ще формулираме и докажем основните свойства на очакването.

\begin{keylem}[Свойства на очакването]
Нека $X$ и $Y$ са две дискретни случайни величини, за които очакванията $\E X$ и $\E Y$ съществуват. Тогава са в сила следните свойства:
\begin{enumerate}
    \item[а)] Ако $X = C$ (константа), то $\E X = C$.
    \item[б)] Ако $c \in \mathbb{R}$ е константа, то $\E(cX) = c\E X$.
    \item[в)] $\E(X + Y) = \E X + \E Y$ (линейност на очакването).
    \item[г)] Ако в допълнение $X$ и $Y$ са \textbf{независими} ($X \perp\!\!\!\perp Y$), то $\E(XY) = \E X \cdot \E Y$.
\end{enumerate}
\end{keylem}

\begin{proof}
а) Ако $X$ приема само една единствена стойност $C$, то $\mathbb{P}(X=C) = 1$. Следователно, по дефиниция:
\[ \E X = C \cdot \mathbb{P}(X=C) = C \cdot 1 = C \]

б) Нека $Z = cX$. Можем да разглеждаме $Z$ като функция на $X$, т.е. $Z = g(X)$, където $g(x) = cx$. Очакването на функция от случайна величина се пресмята като:
\[ \E Z = \sum_{j} g(x_j) p_j = \sum_{j} c x_j p_j = c \sum_{j} x_j p_j = c\E X \]
Константата просто се изнася пред сумата.

в) Нека $Z = X + Y = g(X,Y)$, където функцията е $g(x,y) = x+y$. Използваме дефиницията за математическо очакване за функция от две случайни величини. Обхождаме всички възможни двойки стойности $x_i$ и $y_j$:
\[ \E(X+Y) = \sum_{i, j} (x_i + y_j) \mathbb{P}(X=x_i, Y=y_j) \]
Разкриваме скобите:
\[ = \sum_{i, j} x_i \mathbb{P}(X=x_i, Y=y_j) + \sum_{i, j} y_j \mathbb{P}(X=x_i, Y=y_j) \]
Сега ще прегрупираме сумите. Разглеждаме първата сума и първо сумираме по $j$:
\[ = \sum_{i} x_i \left( \sum_{j} \mathbb{P}(X=x_i, Y=y_j) \right) + \sum_{j} y_j \left( \sum_{i} \mathbb{P}(X=x_i, Y=y_j) \right) \]
Нека обърнем внимание на вътрешната сума $\sum_{j} \mathbb{P}(X=x_i, Y=y_j)$. Тъй като събитията $\{Y=y_j\}$ за всички възможни $j$ образуват пълна група от събития (те разделят достоверното събитие $\Omega$ на непресичащи се части), то обединението им е $\Omega$. Следователно събитието $\{X=x_i\}$ може да бъде записано като обединение на сеченията $\{X=x_i\} \cap \{Y=y_j\}$. Тогава:
\[ \sum_{j} \mathbb{P}(X=x_i, Y=y_j) = \mathbb{P}(X=x_i) \]
Прилагайки същото логическо разсъждение и за втората сума, получаваме:
\[ = \sum_{i} x_i \mathbb{P}(X=x_i) + \sum_{j} y_j \mathbb{P}(Y=y_j) = \E X + \E Y \]

г) За независимост. Нека разгледаме $Z = X \cdot Y$.
\[ \E(XY) = \sum_{i, j} x_i y_j \mathbb{P}(X=x_i, Y=y_j) \]
Тъй като по условие $X$ и $Y$ са независими, вероятността на сечението е равна на произведението от вероятностите:
\[ \mathbb{P}(X=x_i, Y=y_j) = \mathbb{P}(X=x_i) \mathbb{P}(Y=y_j) \]
Заместваме в сумата и групираме членовете:
\[ \E(XY) = \sum_{i, j} x_i y_j \mathbb{P}(X=x_i)\mathbb{P}(Y=y_j) = \left( \sum_{i} x_i \mathbb{P}(X=x_i) \right) \left( \sum_{j} y_j \mathbb{P}(Y=y_j) \right) = \E X \cdot \E Y \]

\end{proof}

\begin{prop}[Неотрицателност на очакването]\label{prop:nonneg-exp}
Ако $X$ е неотрицателна случайна величина (тоест вероятността $X \ge 0$ е единица), то $\E X \ge 0$.
Нещо повече, ако $X \ge 0$ и $\E X = 0$, то $\mathbb{P}(X = 0) = 1$.
\end{prop}

\begin{proof}
По дефиниция $\E X = \sum_{j} x_j p_j$. Тъй като всички стойности $x_j$ са неотрицателни и вероятностите $p_j$ са неотрицателни, то цялата сума също е неотрицателна.

За втората част: сумата $\sum_j x_j p_j$ е сума от неотрицателни събираеми и е равна на нула
само ако всяко събираемо е нула. Значи за всяко $j$ с $x_j > 0$ задължително $p_j = 0$;
следователно цялата вероятностна маса е съсредоточена в стойността $0$.
\end{proof}

\begin{supp}[Очакване чрез опашките]\label{supp:tail-sum}
Ако $X$ приема стойности в $\{0, 1, 2, \dots\}$, то
\[ \E X = \sum_{n \ge 0} \mathbb{P}(X > n) . \]
Наистина, $\mathbb{P}(X > n) = \sum_{k > n} \mathbb{P}(X = k)$, и като сменим реда на
сумиране,
\[
\sum_{n \ge 0} \sum_{k > n} \mathbb{P}(X = k)
 = \sum_{k \ge 1} \sum_{n=0}^{k-1} \mathbb{P}(X = k)
 = \sum_{k \ge 1} k\, \mathbb{P}(X = k) = \E X .
\]

\begin{lecfigure}[Смяната на реда на сумиране, нагледно. В стъпаловидната таблица е
разположена по една клетка за всяка двойка $(n,k)$ с $k > n$; в клетката стои
$p_k = \mathbb{P}(X=k)$. Съберем ли \emph{по редове}, ред $n$ дава опашката
$\mathbb{P}(X > n)$ — това е лявата страна. Съберем ли \emph{по стълбове}, стълб $k$ съдържа
точно $k$ клетки (за $n = 0, \dots, k-1$) и дава $k\,p_k$ — това е дефиницията на $\E X$.
Едни и същи числа, преброени по два начина.\label{fig:tail-sum}]
\begin{tikzpicture}[x=1.13cm, y=0.66cm]
  \def\NC{5}
  % cells: one for every pair (n,k) with k > n
  \foreach \n in {0,...,4} {
    \pgfmathtruncatemacro{\kstart}{\n+1}
    \foreach \k in {\kstart,...,\NC} {
      \fill[defback] ({\k-1.45},{-\n-0.4}) rectangle ({\k-0.55},{-\n+0.4});
      \draw[defframe, line width=0.4pt] ({\k-1.45},{-\n-0.4}) rectangle ({\k-0.55},{-\n+0.4});
      \node[font=\scriptsize] at ({\k-1},{-\n}) {$p_{\k}$};}
    % row label = the tail
    \node[left, font=\scriptsize] at (-0.62,{-\n}) {$\mathbb{P}(X > \n) =$};
    % the row continues
    \node[font=\scriptsize] at ({\NC+0.15},{-\n}) {$\cdots$};}
  \node[left, font=\scriptsize] at (-0.62,-5.05) {$\vdots$};
  \foreach \k in {1,...,\NC} { \node[font=\scriptsize] at ({\k-1},-5.05) {$\vdots$}; }
  % column totals
  \draw[exrule, line width=0.6pt] (-0.5,-5.5) -- ({\NC+0.45},-5.5);
  \foreach \k in {1,...,\NC} {
    \node[font=\scriptsize, below] at ({\k-1},-5.62) {$\k\,p_{\k}$}; }
  \node[font=\scriptsize, below] at ({\NC+0.15},-5.62) {$\cdots$};
  \node[left, font=\scriptsize] at (-0.62,-6.05) {стълбови суми:};
  % the two readings
  \node[right, font=\small, defframe, align=left] at ({\NC+0.75},-1.4)
       {по редове:\\[1pt] $\sum_{n \ge 0} \mathbb{P}(X > n)$};
  \node[right, font=\small, defframe, align=left] at ({\NC+0.75},-4.3)
       {по стълбове:\\[1pt] $\sum_{k \ge 1} k\,p_k = \E X$};
\end{tikzpicture}
\end{lecfigure}
Формулата често спестява работа. Например при $X \sim \Geo(p)$ имаме
$\mathbb{P}(X > n) = q^{\,n+1}$, откъдето веднага
$\E X = \sum_{n \ge 0} q^{\,n+1} = \frac{q}{1-q} = \frac{q}{p}$,
без да се налага да сумираме реда $\sum_k k q^k p$.
\end{supp}

\section{Пример от практиката: Групово тестване (Pool Testing)}

Този пример беше особено актуален по време на пандемията от COVID-19. Представете си, че искаме да тестваме популация, в която хората са здрави (незаразени) с вероятност $p = 0{,}95$ и са заразени с вероятност $1-p = 0{,}05$. Допускаме, че статусите на отделните хора са независими един от друг (което е разумно, освен ако не идват от едно и също домакинство).

Вместо да тестваме всеки човек индивидуално, можем да вземем проби от група от $n$ души и да ги смесим в един общ тест.
\begin{itemize}
    \item Ако общият тест е \emph{отрицателен}, това означава, че всички $n$ човека са здрави. Поради независимостта, вероятността за това събитие е $p^n$. В този случай сме изразходвали само 1 тест за всичките $n$ души.
    \item Ако общият тест е \emph{положителен}, което се случва с вероятност $1 - p^n$, това означава, че поне един човек от групата е заразен. За да разберем кой точно е заразен, се налага да тестваме всеки от тях поотделно. Така сме изразходвали 1 първоначален тест плюс още $n$ индивидуални теста, което прави общо $n+1$ теста.
\end{itemize}

Нека $X_n$ е случайната величина, отчитаща броя на изразходваните тестове за група от $n$ човека. Тя има следното разпределение:
\[
X_n =
\begin{cases}
1, & \text{с вероятност } p^n \\
n+1, & \text{с вероятност } 1-p^n
\end{cases}
\]

Интересува ни средният (очакваният) брой тестове, който се пада на един човек. Това се дава от отношението $\frac{\E X_n}{n}$:
\[ \frac{\E X_n}{n} = \frac{1}{n} \Big( 1 \cdot p^n + (n+1)(1-p^n) \Big) \]
Нека опростим този израз:
\[ \frac{1}{n} \Big( p^n + n + 1 - np^n - p^n \Big) = \frac{n + 1 - np^n}{n} = 1 + \frac{1}{n} - p^n \]

Следователно, пропорцията изразходвани тестове спрямо броя на хората е $1 + \frac{1}{n} - p^n$. Нашата цел е да \textbf{минимизираме} този израз като изберем оптималния размер на групата $n$ (при $n \ge 1$).
Ако разгледаме функцията $f(x) = 1 + \frac{1}{x} - p^x$ и я изследваме за минимум чрез нейната производна ($f'(x) = 0$), за стойност $p = 0{,}95$ ще установим, че минимумът се достига най-близо до цялото число $n = 5$.

Нека пресметнем средния брой тестове на човек при $n = 5$:
\[ \min_{n \ge 1} \frac{\E X_n}{n} = \frac{\E X_5}{5} \approx 0{,}43 \]
Това означава, че на един човек се падат средно $0{,}43$ теста. Оттук намираме очаквания брой тестове за група от петима:
\[ \E X_5 = 5 \cdot 0{,}43 = 2{,}15 \]\footnote{Числото $0{,}43$ е закръглено, затова и произведението е закръглено нагоре. Точната стойност е $\E X_5 = p^5 + 6(1-p^5) = 6 - 5 \cdot 0{,}95^5 \approx 2{,}131$. Изводът не се променя.}

Извод: Вместо да тестваме 5 човека с 5 теста, използвайки тази стратегия, ние очакваме да изразходваме средно само $2{,}15$ теста. Ако имаме 5000 души, ние бихме ги тествали със средно около 2150 теста (вместо 5000). Това е огромно спестяване на ресурси, напълно оправдано чисто математически.

\section{Дисперсия}

Както видяхме в пример~\ref{ex:05-3} с рулетката, математическото очакване само по себе си не винаги е достатъчно, за да опише пълноценно поведението на една случайна величина. То показва къде се намира „центърът на тежестта“, но не ни казва колко са „разпръснати“ възможните стойности около този център. За целта въвеждаме понятието \textbf{дисперсия}.

\subsection{Дефиниция на дисперсия}

\begin{keydefn}
Нека $X$ е дискретна случайна величина, която притежава крайно математическо очакване $\E X$. Тогава \emph{дисперсията} на $X$, която обозначаваме с $\Var X$, се дефинира като математическото очакване на квадрата на отклонението на $X$ от нейното очакване:
\[ \Var X := \E(X - \E X)^2 = \sum_{j} (x_j - \E X)^2 p_j \]
стига тази сума да е крайна, т.е. $\Var X < \infty$.
\end{keydefn}

Ако $X$ приема краен брой стойности, дисперсията винаги съществува. Дисперсията на практика представлява минимизираната средноквадратична грешка (за която говорихме по-рано), когато за характеризираща константа изберем точно математическото очакване $a = \E X$.

\begin{defn}[стандартно отклонение]\label{def:stddev}
Корен квадратен от дисперсията се нарича \emph{стандартно отклонение} (standard deviation):
\[ \sigma_X = \sqrt{\Var X} . \]
То има същата мерна единица като самата случайна величина.
\end{defn}

\subsection{Интерпретация на дисперсията}

Нека се върнем към пример~\ref{ex:05-3} с рулетката:
\begin{itemize}
    \item За първата стратегия ($X$, залагане на цвят): стойностите са $-1$ и $1$. Очакването е $-1/37$. Отклонението на възможните печалби/загуби от очакването е много малко. Следователно $\Var X$ ще е относително малка (около 1).
    \item За втората стратегия ($Y$, залагане на число): стойностите са $-1$ и $35$. Очакването отново е $-1/37$. Въпреки че средната загуба е същата, тук размахът на възможните стойности е огромен. Ако извадите очакването от 35, получавате голямо отклонение, което повдигнато на квадрат дава още по-голям принос към дисперсията. Следователно $\Var Y \gg \Var X$.
\end{itemize}
Дисперсията отчита \emph{волатилността}. В първата игра вие бавно и сигурно губите парите си с малки стъпки. Във втората игра имате шанса да реализирате голяма печалба, макар че дългосрочното математическо очакване остава същото. Разликата в поведението ще стане още по-ясна, когато се изучава Централната гранична теорема.

\begin{example}[Случайно блуждаене (Random Walk)]\label{ex:05-4}
Представете си център на едно въже и $n$ на брой човека, всеки от които може да дърпа въжето с единична сила. Всеки човек се позиционира случайно: отдясно на центъра с вероятност $1/2$ и отляво с вероятност $1/2$.
Нека $X_j$ е силата на $j$-тия човек. Ако е отдясно, $X_j = 1$; ако е отляво, $X_j = -1$.
Резултантната сила е сумата:
\[ F = \sum_{j=1}^n X_j \]
Математическото очакване на силата на един човек е $\E X_j = 1 \cdot \frac{1}{2} + (-1) \cdot \frac{1}{2} = 0$. Поради линейността на очакването, средната резултантна сила също е нула: $\E F = 0$.
Някои хора погрешно смятат, че щом средната сила е нула, то с нарастването на броя на хората $n$, вероятността силата да е точно нула клони към 1. Това е груба грешка. Напротив, дисперсията на силата е много голяма и расте линейно с $n$ ($\Var F = n \cdot \Var X_1 = n$). Следователно стандартното отклонение расте като $\sqrt{n}$. Вероятността частицата (въжето) да остане в покой всъщност е от порядъка на $1/\sqrt{n}$, което клони към $0$ при големи $n$. Частицата прави силно хаотично движение (блуждаене) напред-назад, без да стои на едно място.

\end{example}

\subsection{Свойства на дисперсията}

Ще докажем няколко важни свойства, които значително улесняват работата с дисперсии.

\begin{keythm}[Алтернативна формула за пресмятане на дисперсия]
Ако $X$ е дискретна случайна величина с крайна дисперсия (т.е. $\E X < \infty$ и $\Var X < \infty$), тогава:
\[ \Var X = \E X^2 - (\E X)^2 \]
\end{keythm}

\begin{proof}
Изхождаме от дефиницията и разкриваме скобите под очакването:
\[ \Var X = \E(X - \E X)^2 = \E\big( X^2 - 2X\E X + (\E X)^2 \big) \]
Използваме факта, че $\E X$ е просто константа, и прилагаме линейността на математическото очакване:
\[ = \E X^2 - \E(2X\E X) + \E\big( (\E X)^2 \big) \]
\[ = \E X^2 - 2\E X \cdot \E X + (\E X)^2 \]
\[ = \E X^2 - 2(\E X)^2 + (\E X)^2 = \E X^2 - (\E X)^2 \]
\end{proof}

Тази формула често е много по-удобна за практически пресмятания от самата дефиниция.

\begin{cor}
Нека $X$ е дискретна случайна величина и $\Var X < \infty$. Тогава:
\[ \E X^2 \ge (\E X)^2 \]
\end{cor}

\begin{proof}
Тъй като величината $Y = (X - \E X)^2$ е винаги неотрицателна, нейното очакване също е неотрицателно. Следователно $0 \le \Var X = \E Y = \E X^2 - (\E X)^2$. Като прехвърлим от другата страна, получаваме желаното неравенство.
\end{proof}

\begin{prop}[Свойства на дисперсията]\label{prop:var-props}
Нека $X$ и $Y$ са дискретни случайни величини с крайни дисперсии. Тогава са в сила следните свойства:
\begin{enumerate}
    \item[а)] $\Var X = 0$ тогава и само тогава, когато $X$ е константа с вероятност
    единица, т.е. $\mathbb{P}(X = \E X) = 1$.
    \item[б)] Ако $Z = cX$ за някаква константа $c \in \mathbb{R}$, то $\Var Z = c^2 \Var X$.
    \item[в)] Ако в допълнение $X$ и $Y$ са \textbf{независими} ($X \perp\!\!\!\perp Y$), то $\Var(X + Y) = \Var X + \Var Y$.
\end{enumerate}
\end{prop}

\begin{proof}
а) Ако $X = C$, знаем, че $\E X = C$. Замествайки в дефиницията за дисперсия, получаваме:
\[ \Var X = \E(C - C)^2 = \E(0) = 0 \]
Това е естествено — щом величината е константа, тя не се отклонява от средната си стойност, следователно разсейването (дисперсията) е нулево.
Обратно, ако $\Var X = 0$, то величината $Y = (X - \E X)^2$ е неотрицателна с нулево
очакване, откъдето по твърдение~\ref{prop:nonneg-exp} следва $\mathbb{P}(Y = 0) = 1$, тоест
$X = \E X$ с вероятност единица.

б) За $Z = cX$, намираме дисперсията:
\[ \Var(cX) = \E(cX)^2 - (\E(cX))^2 \]
\[ = \E(c^2 X^2) - (c\E X)^2 \]
\[ = c^2 \E X^2 - c^2 (\E X)^2 = c^2 \Big( \E X^2 - (\E X)^2 \Big) = c^2 \Var X \]
Важно е да се отбележи, че константата излиза от дисперсията повдигната на квадрат. Това е логично, тъй като мерната единица на дисперсията е на втора степен спрямо оригиналната величина (ако мерите в метри, дисперсията е в квадратни метри).

в) Да намерим дисперсията на сума от две независими случайни величини $X$ и $Y$:
\[ \Var(X + Y) = \E(X + Y)^2 - \big(\E(X + Y)\big)^2 \]
Разкриваме скобите:
\[ = \E(X^2 + 2XY + Y^2) - (\E X + \E Y)^2 \]
Използваме линейността на очакването:
\[ = \E X^2 + 2\E(XY) + \E Y^2 - \Big( (\E X)^2 + 2\E X\E Y + (\E Y)^2 \Big) \]
Групираме членовете:
\[ = \Big( \E X^2 - (\E X)^2 \Big) + \Big( \E Y^2 - (\E Y)^2 \Big) + 2\E(XY) - 2\E X\E Y \]
Първите две групи са точно $\Var X$ и $\Var Y$. Остава смесеният член: $2\big(\E(XY) - \E X\E Y\big)$.
Тъй като по условие $X$ и $Y$ са \textbf{независими}, знаем от свойствата на очакването, че $\E(XY) = \E X \cdot \E Y$. Следователно смесеният член се нулира. Така окончателно получаваме:
\[ \Var(X + Y) = \Var X + \Var Y \]
\end{proof}

Това свойство ни показва, че \emph{дисперсията е адитивна, но само за независими случайни величини}. Може да се направи геометрична аналогия: независимостта играе ролята на „ортогоналност“ на вектори. За ортогонални вектори квадратът от дължината на тяхната сума е равен на сумата от квадратите на техните дължини (подобно на Питагоровата теорема). Дисперсията е еквивалентът на този квадрат на дължината.

\section{Задачи}

В модела за групово тестване пренебрегнете биологичните ограничения и за $p=0{,}95$ намерете целия брой хора $n\geq 1$, който минимизира средния брой тестове на човек,
\[
1+\frac{1}{n}-p^n.
\]
Решете задачата числено, например чрез графика. Като алтернативен подход минимизирайте непрекъснатото продължение за $x\geq 1$ и сравнете стойностите при целите числа около получения минимум.
